% !TEX program = xelatex
% Generated by docs/latex/build.py from the sources pinned in sources.json.
\documentclass[11pt,a4paper]{article}
\usepackage[a4paper,margin=25mm]{geometry}
\usepackage{fontspec}
\setmainfont{lmroman10-regular.otf}[BoldFont=lmroman10-bold.otf,ItalicFont=lmroman10-italic.otf,BoldItalicFont=lmroman10-bolditalic.otf]
\setsansfont{lmsans10-regular.otf}[BoldFont=lmsans10-bold.otf,ItalicFont=lmsans10-oblique.otf]
\setmonofont{JuliaMono-Regular.ttf}[Path=../fonts/,Scale=0.83]
\usepackage{amsmath,amssymb,mathtools}
\usepackage{graphicx,calc,longtable,booktabs,array}
\usepackage{fvextra}
\usepackage{xurl}
\usepackage[colorlinks=true,linkcolor=blue,urlcolor=blue,citecolor=blue]{hyperref}
\providecommand{\tightlist}{\setlength{\itemsep}{0pt}\setlength{\parskip}{0pt}}
\setlength{\parindent}{0pt}
\setlength{\parskip}{6pt}
\setlength{\emergencystretch}{6em}
\begin{document}
\section{P3 C5 --- Two cards above a triangular number}\label{p3-c5-two-cards-above-a-triangular-number}

\textbf{DRAFT --- NOT A FULL SOLUTION. Do not submit as a solved C5 answer.} The published work establishes exact finite cases and a general attaining lower-bound family. It does not prove the matching upper bound for all \(k\ge7\) or finish the requested comparison with the C3 method.

Write \(T_j=j(j+1)/2\) and let \(D_B(n)\) be the maximum, over partitions of \(n\), of the number of Bulgarian-solitaire moves before the first cyclic partition. The C5 target recorded in the \href{https://github.com/mpelteshki/climbing-to-the-frontier/blob/212bba53b1818c08c40b8c4b48f636651eaf6410/cagent/p3/c5/README.md}{published package} concerns \(n=T_{k-1}+2\). A move removes one card from each pile, makes a new pile from the removed cards, discards empty piles, and sorts. The exact live C5 page was unavailable during this draft's audit; its wording must be checked before any final submission.

\subsection{Exact finite results}\label{exact-finite-results}

The six \href{https://github.com/mpelteshki/climbing-to-the-frontier/blob/212bba53b1818c08c40b8c4b48f636651eaf6410/cagent/p3/README.md}{Lean theorems and replay instructions} establish the following exact maxima over \textbf{all} partitions at the listed sizes. The \href{https://github.com/mpelteshki/climbing-to-the-frontier/blob/212bba53b1818c08c40b8c4b48f636651eaf6410/cagent/p3/evidence/verification.json}{recorded verification log} reports successful checks. These are six individual statements, not an induction over \(k\).

{\def\LTcaptype{none} % do not increment counter
\begin{longtable}[]{@{}
  >{\raggedright\arraybackslash}p{(\linewidth - 6\tabcolsep) * \real{0.2143}}
  >{\raggedleft\arraybackslash}p{(\linewidth - 6\tabcolsep) * \real{0.2857}}
  >{\raggedleft\arraybackslash}p{(\linewidth - 6\tabcolsep) * \real{0.2857}}
  >{\raggedright\arraybackslash}p{(\linewidth - 6\tabcolsep) * \real{0.2143}}@{}}
\toprule\noalign{}
\begin{minipage}[b]{\linewidth}\raggedright
\(k\)
\end{minipage} & \begin{minipage}[b]{\linewidth}\raggedleft
\(n=T_{k-1}+2\)
\end{minipage} & \begin{minipage}[b]{\linewidth}\raggedleft
\(D_B(n)\)
\end{minipage} & \begin{minipage}[b]{\linewidth}\raggedright
Lean theorem
\end{minipage} \\
\midrule\noalign{}
\endhead
\bottomrule\noalign{}
\endlastfoot
\(2\) & \(3\) & \(2\) & \href{https://github.com/mpelteshki/climbing-to-the-frontier/blob/212bba53b1818c08c40b8c4b48f636651eaf6410/cagent/p3/lean/ProofPursuit/P3/N3.lean}{\texttt{maximumDepth\_3}} \\
\(3\) & \(5\) & \(3\) & \href{https://github.com/mpelteshki/climbing-to-the-frontier/blob/212bba53b1818c08c40b8c4b48f636651eaf6410/cagent/p3/lean/ProofPursuit/P3/N5.lean}{\texttt{maximumDepth\_5}} \\
\(4\) & \(8\) & \(5\) & \href{https://github.com/mpelteshki/climbing-to-the-frontier/blob/212bba53b1818c08c40b8c4b48f636651eaf6410/cagent/p3/lean/ProofPursuit/P3/N8.lean}{\texttt{maximumDepth\_8}} \\
\(5\) & \(12\) & \(8\) & \href{https://github.com/mpelteshki/climbing-to-the-frontier/blob/212bba53b1818c08c40b8c4b48f636651eaf6410/cagent/p3/lean/ProofPursuit/P3/N12.lean}{\texttt{maximumDepth\_12}} \\
\(6\) & \(17\) & \(12\) & \href{https://github.com/mpelteshki/climbing-to-the-frontier/blob/212bba53b1818c08c40b8c4b48f636651eaf6410/cagent/p3/lean/ProofPursuit/P3/N17.lean}{\texttt{maximumDepth\_17}} \\
\(7\) & \(23\) & \(18\) & \href{https://github.com/mpelteshki/climbing-to-the-frontier/blob/212bba53b1818c08c40b8c4b48f636651eaf6410/cagent/p3/lean/ProofPursuit/P3/N23.lean}{\texttt{maximumDepth\_23}} \\
\end{longtable}
}

\subsection{General lower-bound witness}\label{general-lower-bound-witness}

The following is the \href{https://github.com/mpelteshki/climbing-to-the-frontier/blob/212bba53b1818c08c40b8c4b48f636651eaf6410/cagent/p3/c5/lower-bound-proof.md}{published written witness argument}, with its literature attribution narrowed to the range covered by the cited case. It proves \(D_B(T_{k-1}+2)\ge(k-1)(k-4)\) for every \(k\ge5\).

For every \(k\ge5\), the partition

\[
\lambda_k=(k-2,k-2,k-3,\ldots,3,3,2,1)
\]

has first cyclic depth exactly \((k-1)(k-4)\). To remove ambiguity in the short tail, its \(k\) entries are \(\lambda_1=k-2\), \(\lambda_i=k-i\) for \(2\le i\le k-3\), and \((\lambda_{k-2},\lambda_{k-1},\lambda_k)=(3,2,1)\). Their sum is \(T_{k-1}+2\). The family agrees with the \(r=2\) specialization of \href{https://sc.edu/study/colleges_schools/artsandsciences/mathematics/research/imi/research/documents/1998/1998_12.pdf}{Griggs and Ho, Theorem 4.5, Case (1)} for \(k\ge6\); the direct calculation below also covers \(k=5\). For \(k=5,6\) this family is not optimal.

Use columns numbered from zero and rows numbered from one. Relative to the staircase of height \(k-1\), the initial diagram has one hole on diagonal \(k-1\) at column \(0\), and three extra cards on diagonal \(k\) at columns \(k-3,k-2,k-1\).

Let \(F=(k-1)(k-4)\). Before sorting, a Bulgarian move rotates diagonal \(w\) one column forward modulo \(w\). At a time \(0\le t<F\), write \(t=a(k-1)+b\) with \(0\le a\le k-5\) and \(0\le b\le k-2\). The hole rotates to column \(b\). The extra cards rotate to

\[
(b-a-3)\bmod k,\quad(b-a-2)\bmod k,\quad(b-a-1)\bmod k.
\]

Each unreduced value is either nonnegative and strictly less than \(b\), or negative and becomes at least \(b+2\) after adding \(k\): the smallest wrapped difference is \(k-a-3\ge2\). Thus none of the extras occupies the hole's column or the column immediately to its right. There is no vertical gap, and the column heights remain decreasing. Away from the hole, a staircase with distinct extra columns cannot have an adjacent inversion; at the hole, only an extra immediately to its right could cause one. Sorting therefore does nothing before time \(F\), so induction proves this trajectory formula at every such time.

At \(t=F\), the unsorted diagram has the hole at column \(0\) and extra cards at columns \(1,2,3\). The first four heights are \(k-2,k-1,k-2,k-3\). Sorting swaps the first two, giving

\[
(k-1,k-2,k-2,k-3,k-5,k-6,\ldots,1).
\]

The tail is empty when \(k=5\). This is a complete staircase of height \(k-1\) plus one card in each of columns \(2,3\), hence a cyclic partition by the \href{https://github.com/mpelteshki/climbing-to-the-frontier/blob/212bba53b1818c08c40b8c4b48f636651eaf6410/cagent/p3/c1/written-proof.md}{C1 boundary classification}. Before time \(F\), a hole remained below that boundary, so no earlier state was cyclic. The witness therefore has first cyclic depth exactly \(F\).

\subsection{Unresolved obligations}\label{unresolved-obligations}

A complete C5 solution still needs a proof that \textbf{every} partition of \(T_{k-1}+2\) has depth at most \((k-1)(k-4)\) for every \(k\ge7\). The \(k=7\) Lean instance and additional finite checks cannot prove that universal bound. The \href{https://github.com/mpelteshki/climbing-to-the-frontier/blob/212bba53b1818c08c40b8c4b48f636651eaf6410/cagent/p3/c3/written-proof.md}{C3 general nontriangular bound} gives only \(D_B(n)\le k^2-2k-1\) at rank \(k\), which is \(3k-5\) above the proposed C5 value \(k^2-5k+4\). This arithmetic comparison does not explain a valid refinement of C3's pile-count method for \(r=2\); that requested methodological comparison is also open.

Griggs and Ho state their general expression as a \textbf{lower bound} in Theorem 4.5 and equality as \href{https://sc.edu/study/colleges_schools/artsandsciences/mathematics/research/imi/research/documents/1998/1998_12.pdf}{Conjecture 4.7}. The cited paper therefore supplies no missing upper-bound proof. \textbf{C5 remains Not solved.}
\end{document}
